\section{Two quadratic cube descents and exact elliptic ranks}
\label{gd-section}

Keep the elliptic curves and three-isogeny of
Theorem~\ref{ssg-isogeny}:
\[
 E:Y^2=f(X)=X^3-9X-9,\qquad
 E':V^2=f_*(T)=T^3-189T+999,
\]
\[
 \phi(X,Y)=(F(X),YF'(X)),\qquad
 F(X)=X+\frac{36}{X+3}-\frac{36}{(X+3)^2}.
\]
Put \(\epsilon=2+\sqrt3\). We prove exact rank statements by explicit
descent. These statements do not determine the rational points of the
simultaneous quartic in Theorem~\ref{ssg-quartic}.

\begin{theorem}[The automatic Gaussian cube]\label{gd-gaussian}
For every finite \((X,Y)\in E(\mathbb Q)\), there is a unique
\(\alpha+i\beta\in\mathbb Q(i)\) such that
\begin{equation}\label{gd-gaussian-cube}
 (\alpha+i\beta)^3=Y+3i(X+2).
\end{equation}
\end{theorem}
\begin{proof}
Use primitive coordinates
\[
 X=A/d^2,\quad Y=B/d^3,\quad d>0,\quad
 \gcd(A,d)=\gcd(B,d)=1,\quad
 B^2=A^3-9Ad^4-9d^6.
\]
One has \(3\nmid B\). Otherwise \(3\mid A\), so \(3\nmid d\), and
reduction modulo \(27\) gives \(B^2\equiv18\pmod{27}\). Writing
\(B=3b\) would give \(b^2\equiv2\pmod3\), a contradiction.
Set \(K=A+2d^2\) and \(S=B+3idK\). Then
\begin{equation}\label{gd-gaussian-integral}
 N(S)=B^2+9d^2K^2=(A+3d^2)^3>0.
\end{equation}
For a prime \(q>3\) dividing \(dK\), one has \(q\nmid B\): if
\(q\mid d\), the curve equation gives \(B^2\equiv A^3\), a unit;
if \(q\mid K\), then \(d\) is a unit and \(B^2\equiv d^6\).
Thus no such \(q\) divides both coefficients of \(S\).

In \(\mathbb Z[i]\), the unique prime above \(2\) has norm \(2\);
its exponent in \(S\) is \(v_2(N(S))\), hence a multiple of three.
There is no prime above \(3\) in \(S\). For \(q>3\), an inert prime
cannot divide \(S\), since this would divide both coefficients.
At a split prime at most one conjugate orientation occurs, for the
same reason; its exponent is \(v_q(N(S))\), again divisible by three.
These quadratic splitting and ramification facts follow from
\cite{MilneANT2020}, Examples 2.41 and 3.44--3.45.

The ring is norm-Euclidean: rounding both coordinates leaves a
difference of norm at most \(1/2\). It is therefore a UFD.
Its four units are permuted by cubing. Consequently \(S=W^3\)
for an actual Gaussian integer \(W\), and \(W/d\) proves existence
in \eqref{gd-gaussian-cube}. This root is nonzero. The field
\(\mathbb Q(i)\) has no nontrivial cube root of unity, so the root
is unique.
\end{proof}
The cube condition is automatic on \(E(\mathbb Q)\); it is not an
additional exclusion of elliptic points.

\begin{theorem}[The opposite isogeny is bijective on rational points]
\label{gd-opposite}
Define
\begin{equation}\label{gd-psi}
 G(T)=\frac{T+108/(T-9)+108/(T-9)^2}{9},\qquad
 c(T)=\frac{1-108/(T-9)^2-216/(T-9)^3}{27}.
\end{equation}
The map \(\psi(T,V)=(G(T),Vc(T))\) extends to a degree-three
isogeny \(E'\to E\), induces a bijection
\(E'(\mathbb Q)\to E(\mathbb Q)\), and satisfies
\(\psi\phi=[3]\).
For a finite point on \(E\), the inverse is
\begin{equation}\label{gd-psi-inverse}
 T=9+\frac3{\beta-1},\qquad V=\frac{9\alpha}{\beta-1},
\end{equation}
where \(\alpha+i\beta\) is the root in
Theorem~\ref{gd-gaussian}.
\end{theorem}
\begin{proof}
Clearing denominators proves \(f_*(T)c(T)^2=f(G(T))\).
The rational map extends to the smooth projective curves and maps
\(O\) to \(O\), so it is a homomorphism by
\cite{MilneEC2Descent}, Chapter II, Proposition 1.5.
For \(u=T-9\), one has \(G=1+u/9+12/u+12/u^2\).
Its reduced numerator is cubic and nonzero at \(u=0\); the abscissa
degree diagram gives degree three. Its geometric kernel consists of
\(O,(9,3\sqrt3),(9,-3\sqrt3)\), so its rational kernel is trivial.

Norm and imaginary-coordinate comparison in
\eqref{gd-gaussian-cube} give
\[
 \alpha^2+\beta^2=X+3,\qquad
 \beta(3\alpha^2-\beta^2)=3(X+2).
\]
If \(\beta=1\), these imply \(3X+5=3X+6\).
Thus \eqref{gd-psi-inverse} is defined. Substitution proves its
curve equation and its image under \(\psi\).
For explicit checks in the reverse direction, if \(u=T-9\), set
\(\beta=1+3/u\), \(\alpha=V/(3u)\). Then
\begin{align}
 \alpha^2&=\frac{f_*(T)}{9u^2},
 &\alpha^2+\beta^2&=G(T)+3,\label{gd-gaussian-inverse-identities}\\
 \beta(3\alpha^2-\beta^2)&=3(G(T)+2),
 &\alpha(\alpha^2-3\beta^2)&=Vc(T).\notag
\end{align}
Conversely the two norm and imaginary equations imply
\(4\beta^3-3(X+3)\beta+3(X+2)=0\); with \(\beta\ne1\), this
solves for \(X\) as \(G(T)\) in \eqref{gd-psi-inverse}.
The real coordinate recovers \(Y\).
Every finite rational point therefore has a rational preimage.
The only rational preimage of \(O\) is \(O\), and the trivial
rational kernel proves uniqueness everywhere. The model scalings
\(1/9\) and \(1/27\) in \eqref{gd-psi} are retained.

Finally, on a dense chart put
\begin{align*}
 x_2&=\frac{(3X^2-9)^2}{4f(X)}-2X,&
 b_2&=\frac{(3X^2-9)(X-x_2)}{2f(X)}-1,\\
 \ell&=\frac{b_2-1}{x_2-X},&
 x_3&=f(X)\ell^2-X-x_2,\qquad
 b_3=\ell(X-x_3)-1.
\end{align*}
The chord law gives \([3](X,Y)=(x_3,Yb_3)\).
The exact rational-function identities
\begin{equation}\label{gd-composition}
 G(F(X))=x_3,\qquad F'(X)c(F(X))=b_3
\end{equation}
follow by clearing denominators. Equality of morphisms on this dense
chart extends over every omitted point. Thus \(\psi\phi=[3]\),
with the asserted orientation.
\end{proof}

\begin{lemma}[Three real-quadratic classes]\label{gd-three-classes}
For every finite \((T,V)\in E'(\mathbb Q)\), the nonzero element
\[
 \theta(T,V)=V+3\sqrt3(T-8)
\]
has cube class in the three distinct classes
\(1,\epsilon,\epsilon^2\) in
\(\mathbb Q(\sqrt3)^\times/\mathbb Q(\sqrt3)^{\times3}\).
\end{lemma}
\begin{proof}
The value \(T=9\) is impossible because \(V^2=27\).
Choose primitive coordinates \(T=A/d^2,V=B/d^3\).
Then
\begin{equation}\label{gd-real-integral}
 S=B+3\sqrt3d(A-8d^2),\qquad N(S)=(A-9d^2)^3\ne0.
\end{equation}
The integer norm may be negative. For \(q>3\), if \(q\mid d\)
the curve equation gives a unit \(B^2\equiv A^3\); if
\(q\mid A-8d^2\), it gives
\(B^2\equiv(8^3-189\cdot8+999)d^6=-d^6\), a unit.
Thus \(q\) cannot divide both coefficients of \(S\).

The ring of integers is \(\mathbb Z[\sqrt3]\), of discriminant \(12\).
Each of \(2\) and \(3\) has one ramified prime of residue degree one,
so its exponent in \(S\) is a multiple of three by
\eqref{gd-real-integral}. At \(q>3\), inert divisibility is impossible
and only one split orientation occurs; again the exponent is a
multiple of three.

This ring is norm-Euclidean: rounding the two coordinates gives
\(|a|,|b|\leq1/2\), so \(|a^2-3b^2|\leq3/4<1\).
Its units are exactly \(\{\pm\epsilon^j:j\in\mathbb Z\}\).
Indeed norm \(-1\) is impossible modulo \(3\). After changing sign
and dividing by a power of \(\epsilon\), a unit lies in
\([1,\epsilon)\) in the chosen real embedding. A unit strictly
above \(1\) has integral coordinates \(x\geq2,y\geq1\), since its
conjugate is its positive reciprocal; it would therefore be at least
\(\epsilon\). This proves the unit assertion.

Unique factorization now gives \(S=\epsilon^jW^3\), absorbing the
sign into \(W\). Division by \(d^3\) proves the three-class bound.
A cube root of \(\epsilon\) or \(\epsilon^2\) in the field would be
an integral unit by its prime valuations, contradicting the unit
exponents just proved. The three classes are distinct.
\end{proof}

\begin{theorem}[The exact descent quotient]\label{gd-quotient}
With \(\delta(O)=1\) and \(\delta(T,V)=[\theta(T,V)]\), the map
\[
 \delta:E'(\mathbb Q)\longrightarrow
 \mathbb Q(\sqrt3)^\times/\mathbb Q(\sqrt3)^{\times3}
\]
is a homomorphism with kernel \(\phi(E(\mathbb Q))\) and image of
order three. Consequently
\begin{equation}\label{gd-index}
 |E'(\mathbb Q)/\phi(E(\mathbb Q))|=3.
\end{equation}
\end{theorem}
\begin{proof}
Put \(T_0=(9,-3\sqrt3)\). Its tangent is
\(Y=m_0X+b_0\), with \(m_0=-3\sqrt3,b_0=24\sqrt3\);
the line function is \(\theta=Y-m_0X-b_0\).
One has the exact flex identity
\[
 X^3-189X+999-(m_0X+b_0)^2=(X-9)^3.
\]
Let a nonvertical rational chord or tangent be \(Y=mX+b\), with
three intersections \(P,Q,R\), counted with multiplicity, so
\(P+Q+R=O\). Its monic intersection polynomial is
\(F_{\rm line}(X)=X^3-189X+999-(mX+b)^2\).
If \(m\ne m_0\), set \(x_0=-(b-b_0)/(m-m_0)\).
The two lines agree at \(x_0\), so \(F_{\rm line}(x_0)=(x_0-9)^3\).
Multiplication of their linear evaluations at the three roots gives
\begin{equation}\label{gd-chord}
 \theta(P)\theta(Q)\theta(R)
   =[b-b_0+9(m-m_0)]^3.
\end{equation}
For \(m=m_0\), the evaluations are all \(b-b_0\), giving the same
conclusion. This polynomial identity includes repeated intersections.
Its right side is nonzero: a rational line cannot contain \(T_0\),
whose abscissa is rational and ordinate irrational.
Moreover
\begin{equation}\label{gd-inversion}
 \theta(P)\theta(-P)=(9-X(P))^3\ne0.
\end{equation}
Equations \eqref{gd-chord}--\eqref{gd-inversion} prove the
homomorphism; the latter also handles vertical lines, and
\(\delta(O)=1\) handles identity cases.

For the kernel, suppose
\(\theta(T,V)=(u+v\sqrt3)^3\), \(u,v\in\mathbb Q\).
Norm and imaginary-coordinate comparison give
\[
 u^2-3v^2=T-9,\qquad v(u^2+v^2)=T-8.
\]
Equality of rational cubes justifies the first equality for either
sign of the norm. The value \(v=1\) would give both
\(u^2=T-6\) and \(u^2=T-9\). Thus
\begin{equation}\label{gd-phi-inverse}
 X=-3+\frac3{1-v},\qquad Y=\frac{3u}{1-v}
\end{equation}
is defined, and substitution gives a rational point of \(E\)
mapping to \((T,V)\) under \(\phi\).
Explicitly, with \(z=X+3\), \(u=Y/z\), \(v=1-3/z\), the
identities are
\[
 u^2-3v^2=F(X)-9,\quad
 v(u^2+v^2)=F(X)-8,\quad
 u(u^2+9v^2)=YF'(X).
\]
They also prove that every finite image has trivial cube class.
There is no rational point with \(X=-3\), since its ordinate would
satisfy \(Y^2=-9\). The identity point and its preimage \(O\)
complete the kernel statement.

Finally \(P'=(6,9)\) gives
\[
 \theta(P')=9-6\sqrt3=\epsilon^2(3-2\sqrt3)^3.
\]
Its class has order three. Lemma~\ref{gd-three-classes} therefore
gives the exact image and \eqref{gd-index}.
\end{proof}

\begin{corollary}[Exact ranks, without a generator claim]\label{gd-ranks}
The ranks of \(E(\mathbb Q)\) and \(E'(\mathbb Q)\) are both one.
For the nontrivial-unit quotient \(H_1\), the rank of
\(\operatorname{Jac}(H_1)(\mathbb Q)\) is exactly two.
\end{corollary}
\begin{proof}
The bijective rational homomorphism \(\psi\), with \(\psi\phi=[3]\),
identifies \eqref{gd-index} with \(|E(\mathbb Q)/3E(\mathbb Q)|=3\).
By Mordell--Weil (\cite{MilneEC2Descent}, Chapter IV, printed
p.~105), write \(E(\mathbb Q)\simeq\mathbb Z^r\oplus T\) with
\(T\) finite. Then \(3=3^r|T/3T|\). The ordinary EQ doubling proof
already gives the infinite-order point \((-2,1)\), so \(r\geq1\).
It follows that \(r=1\); \(\psi\) gives the same rank for \(E'\).
The established \(\mathbb Q\)-isogeny
\(\operatorname{Jac}(H_1)\sim E\times E\) proves the final assertion.
\end{proof}

Neither this argument nor the finite arithmetic replay proves that
\((-2,1)\) or \((6,9)\) generates its full group. The simultaneous
same-source quartic and its positive inequalities remain to be solved.
The Gaussian cube is a reversible descent, not a new point exclusion.
The exact replay checks eleven rational-function identities, the
displayed unit-class witness, and the complete primitive mod-\(27\)
model table; it uses no software rank upper bound.


\paragraph{The arithmetic formally checked here.}
The module \texttt{GaussianDescent\allowbreak Arithmetic.lean} contains
15 new theorems. A fresh copy of the source was compiled to a fresh
\texttt{olean} with warnings treated as errors, under Lean \(4.32.0\)
and the pinned Mathlib snapshot. All 15 complete axiom queries use
only the standard three axioms.

The module proves the two actual homogeneous integer norm identities,
reduction modulo an arbitrary integer modulus, a complete kernel-checked
\(\operatorname{Fin}(27)\) table, its transfer to signed integer
congruences, and the actual primitive-curve consequence \(3\nmid B\).
It also proves the opposite-isogeny rational-function identity, the
three Gaussian inverse and three real-quadratic inverse identities,
and both coordinate expressions for \(\psi\phi\) and chord-law tripling
on the displayed rational charts. Every denominator assumption remains
explicit in the theorem signatures; no projective exception is hidden
inside a formal global curve assertion.

The quadratic rings, their units and prime splitting, existence and
uniqueness of cube roots, the homomorphism \(\delta\), isogeny degree,
extension of the maps, Mordell--Weil and the exact ranks remain the
reviewed ordinary proof above. In particular these arithmetic theorems
do not constitute a formal proof of the entire isogeny descent or of
the rational points of the simultaneous curve.
